% !PRINTED 494
\weberpage{494}
%<<CONT
are, but that therein, in the case $n = 7$, all the $A_i$ depend rationally on $f(\omega)^{24}$, whereas in the case $n = 11$ the $A_i$ with even index likewise depend rationally on $f(\omega)^{24}$, while those with odd index are the product of $f(\omega)^{12}$ with a rational function of $f(\omega)^{24}$. Apart from rational numbers, the numerical coefficients contain only the irrationality $i\sqrt{7}$ or $i\sqrt{11}$ respectively, and accordingly we will denote these equations by

\weq{17}{\begin{aligned}
\Phi\bigl[w, f(\omega)^{24},\ \ i\sqrt{7}\bigr] &= 0,\\
\Phi\bigl[w, f(\omega)^{12},\ \ i\sqrt{11}\bigr] &= 0.
\end{aligned}}

If we change the signs of $i$ and $w$ at the same time, then, as we have seen in \S\,82, there results an equation whose roots are the quantities $w'$.

In the application of the substitution $(c''')$ also a difference between the two cases shows itself.

For it follows from (13), taking into account that by (2) the product of all the $v_\nu$ has the value $u^8$, that in the case $n = 7$ the $w_\nu$ are merely permuted among one another by the substitution $(c''')$, and indeed in the following manner:

\[
\begin{pmatrix} w_0, & w_1, & w_2, & w_3, & w_4, & w_5, & w_6\\ w_5, & w_3, & w_6, & w_2, & w_4, & w_0, & w_1 \end{pmatrix},
\]

and from this it follows that in this case the coefficients $A_\nu$ all remain unchanged under the interchange:

\[
\left(f(\omega)^{24},\ \frac{2^{12}}{f(\omega)^{24}}\right).
\]

Moreover, in the denominator of these coefficients only powers of $f(\omega)^{24}$ occur (\S\,73), and the powers of $f(\omega)^{24}$ rise to the same height as those of $f(\omega)^{-24}$.

In the case $n = 11$ the quantities $w_\nu$ pass over into the quantities $w'_\nu$ under the permutation $(c''')$, and indeed in the following order:

\[
\begin{pmatrix} w_0, & w_1, & w_2, & w_3, & w_4, & w_5, & w_6, & w_7, & w_8, & w_9, & w_{10}\\ w'_8, & w'_2, & w'_7, & w'_0, & w'_9, & w'_5, & w'_6, & w'_4, & w'_3, & w'_{10}, & w'_1 \end{pmatrix}.
\]

By (17) we can express this property by the equation identical with respect to $w$:

\[
-\Phi\bigl[-w, f(\omega)^{12}, -i\sqrt{11}\bigr] = \Phi\!\left(w, \frac{2^6}{f(\omega)^{12}}, i\sqrt{11}\right),
\]

and from this there result the following properties of the coefficients $A_\nu$:
%CONT>>
